% Generated draft: requires mathematical and source review.
\documentclass[11pt]{article}
\usepackage[margin=25mm]{geometry}
\usepackage{fontspec}
\setmainfont{Latin Modern Roman}
\usepackage{amsmath,amssymb,mathtools,unicode-math}
\setmathfont{Latin Modern Math}
\usepackage{enumitem,needspace,xurl,hyperref}
\setlength{\parindent}{0pt}
\setlength{\parskip}{5pt}
\setlength{\emergencystretch}{3em}
\newcommand{\N}{\mathbb{N}}
\newcommand{\Z}{\mathbb{Z}}
\newcommand{\Q}{\mathbb{Q}}
\newcommand{\R}{\mathbb{R}}
\newcommand{\C}{\mathbb{C}}
\newcommand{\F}{\mathbb{F}}
\newcommand{\A}{\mathbb{A}}
\newcommand{\PP}{\mathbb{P}}
\newcommand{\E}{\mathbb{E}}
\newcommand{\eps}{\varepsilon}
\newcommand{\dd}{\,\mathrm{d}}
\begin{document}
\section*{UnsolvedMath ID: 20002064}
\textbf{Source problem code:} catalogue entry
\subsection{Definitions and mathematical statement}
The source reports Ko Honda's joint results with Etnyre that some knot types are not transversally simple, obtained indirectly, and Eliashberg's request for an invariant distinguishing transverse knots. In a cooriented contact three-manifold, a positively transverse knot has tangent positively transverse to the contact planes. In the standard \((S^3,\xi_{\mathrm{std}})\) setting used by the aid, a smooth oriented knot type is transversally simple if its positive transverse isotopy classes are determined by self-linking number. The general source does not specify an ambient manifold or ask for a complete invariant.

The aid reports that the invariant-existence interpretation has positive answers through Khovanov, Floer, monopole and transverse-DGA constructions, while a complete transverse classification does not follow. It supplies the following abstract descent criterion in a setting where the Legendrian-approximation/negative-stabilization theorem applies. Let \(\mathcal L\) be Legendrian isotopy classes, \(S_-\) negative stabilization, and \(I:\mathcal L\to A\) an invariant satisfying
\[
 I(S_-L)=f(I(L))
\]
for a fixed map \(f:A\to A\). Let \(A/{\sim}\) be the quotient by the smallest equivalence relation identifying \(a\) with \(f(a)\). Positive transverse push-off \(T_+(L)\) then has the well-defined invariant
\[
 \overline I(T_+(L))=[I(L)]\in A/{\sim}.
\]
This uses the theorem that two Legendrian approximations give the same transverse class exactly when some negative stabilizations are Legendrian isotopic. The resulting invariant is complete precisely when
\[
 [I(L)]=[I(L')]
 \quad\Longrightarrow\quad
 S_-^mL\simeq S_-^{m'}L'
 \text{ for some }m,m'\geq0.
\]
The reverse implication is automatic. A quotient construction alone therefore provides descent, not completeness.

For classical Legendrian values \((t,r)=(\operatorname{tb},\operatorname{rot})\), use
\[
 S_-:(t,r)\mapsto(t-1,r-1),\qquad
 \operatorname{sl}(T_+(L))=\operatorname{tb}(L)-\operatorname{rot}(L).
\]
On the formal value lattice \(\Z^2\), the coequalizer is exactly \(t-r\); realizability restricts its image for a given knot type. The aid cites Etnyre--Honda representatives labelled \(L_+\) and \(S_+K_+\) in the \((2,3)\)-cable of the positive trefoil. Their positive transverse push-offs both have self-linking \(3\) but are not negative-stably Legendrian equivalent. Consequently any descended invariant factoring only through classical Legendrian data and the fixed smooth knot type fails to distinguish this pair. The individual representatives and cable convention require the original reference; no unstated \(\operatorname{tb}\) or rotation values are invented. This is an incompleteness certificate, not a claim that every cited modern invariant distinguishes that pair.
\subsection{Short English statement}
The source asks for invariants detecting transverse non-simplicity. The aid reports several existing theories and gives a stabilization-quotient test for descent and completeness. Classical Legendrian data descend only to self-linking, which cannot distinguish the cited Etnyre--Honda pair with self-linking \(3\).
\subsection{Sources}
\begin{itemize}
\item Local UnsolvedMath registry, permanent ID \(20002064\), source code \texttt{catalogue entry}; canonical collection \texttt{aim-geometry-notes.json}, notes[401].
\item AIM workshop ``Holomorphic curves in contact geometry,'' Question \(1.6\): \url{https://aimath.org/WWN/contactgeom2/contactgeom2.pdf}. The source names Ko Honda, Etnyre and Eliashberg. ``Second day'' is the following heading. No PDF was opened in this session.
\item The aid names Khovanov, Floer, monopole and transverse-DGA constructions and the Etnyre--Honda cable representatives, without full authoritative bibliographic entries or theorem statements. Original remarks and literature fields are empty; no proposer-field value or year is supplied.
\item Supplied research attempt \(1\), \texttt{partial\_result}, run status \texttt{partially\_solved}, rigor label ``proved,'' novelty confidence ``low.'' Registry classification \texttt{PARTIAL-PROGRESS}, updated 14 August 2026. The historical ``no invariant here'' description is not a current nonexistence assertion.
\end{itemize}
\end{document}
